\chapter{תקציר מנהלים}\label{ch:exec}

\begin{keybox}[בשורה אחת]
מערכת חדשות שתאמץ את \en{VeriLens} תדע בתוך דקה אם תמונה היא תיעוד אמיתי או זיוף, תוכל להוכיח זאת לקוראים, ותעמוד בחובת הסימון האירופית. הפיתוח נמשך \num{18} חודשים ועולה \num{9.6} מיליון ₪.
\end{keybox}

\begin{notebox}[היקף: חזון מול מטלה]
לפי הנחיית המרצה, בקורס כותבים את הדוח המכין ולא מבצעים את הפרויקט, ומותר ״לחשוב בגדול״. לכן החוזה בדוח הוא חוזה של החברה המתוכננת \en{VeriLens}: \num{18} חודשים, צוות של \num{14} ותקציב של \num{9.6} מיליון ₪. תפקידי ההנהלה (מנכ״ל, \en{CTO}, ראש אתיקה) הם תפקידים בחברה הזאת. מה שאנחנו, גיא ופז, מגישים בקורס הוא הדוח הזה, התבנית, ותהליך העבודה המתועד במאגר.
\end{notebox}

\textbf{מה הלקוח צריך.} עורכי חדשות מקבלים כל יום תמונות ממקורות שאינם מאומתים, ונדרשים להחליט תוך דקות. מחוללי תמונות מייצרים היום ״צילום חדשותי״ של אירוע שלא קרה, והרשת מתגמלת זאת בטראפיק. במאי \num{2023} תמונה מזויפת של פיצוץ ליד הפנטגון הורידה לזמן קצר את מדדי וול סטריט~\cite{npr2023}. \num{125} עמודי פייסבוק צברו מאות מיליוני חשיפות מתמונות סינתטיות~\cite{diresta2024}, וביוני \num{2026} מנתה \en{NewsGuard} כבר \num{3,749} אתרי תוכן אוטומטיים~\cite{newsguard2026}. ללקוח יש שלושה חששות עיקריים: לפרסם זיוף תחת המותג שלו, לאבד את הקוראים, ולהפר את החוק.

\textbf{מה אנחנו מציעים.} שני רכיבים שנמכרים יחד:
\begin{enumerate}
\item \textbf{פלטפורמת אימות} שמשתלבת במערכת ניהול התוכן. כל תמונה עוברת שלוש בדיקות: תעודת מקור (\en{C2PA}), ארבעה גלאים סטטיסטיים, ובדיקת הקשר. העורך מקבל ציון והמלצה אחת מתוך שלוש: ״מאומת״, ״לא ודאי״ או ״חשד גבוה״.
\item \textbf{התקן האתי \en{TRUTH-IMG}}: עשרה כללי עבודה, תיוג גלוי לקוראים, יומן החלטות, והסמכה שנתית.
\end{enumerate}

\textbf{למה דווקא אנחנו.} זיהוי תמונות למערכות חדשות כבר קיים בשוק: \en{ImageWhisperer} מריץ \num{44} בדיקות, בודק גם תעודות \en{C2PA} ונמכר לקבוצות מדיה אירופיות~\cite{imagewhisperer,iwhome}, ו-\en{AP} מוכרת למערכות חדשות את \en{AP Verify}~\cite{apverify}. לכן אנחנו לא מתחרים על ״זיהוי טוב יותר״. טבלה~\ref{tab:discrim} מציגה את שלושת היתרונות שכן מבדילים את ההצעה, את התועלת ללקוח, ואת ההוכחה לכל אחד מהם.

\begin{table}[!htb]\centering
\caption{יתרונות מבדלים: מה מקבל הלקוח ועל מה זה נשען}\label{tab:discrim}
\small
\begin{tabularx}{\linewidth}{@{}Y Y Y@{}}
\toprule
\textbf{מה אנחנו עושים אחרת} & \textbf{מה הלקוח מרוויח} & \textbf{ההוכחה} \\
\midrule
תשובה ברורה לקורא, לא רק ציון לעורך & כל תמונה שפורסמה מקבלת אחת מארבע תוויות, עם תעודת המקור ורשומת ההחלטה בלחיצה & קוראים חושבים ב״אמיתי או לא״, וציון הסתברות קשה להם להבנה~\cite{citp2026}; למתחרים אין תווית לקורא (טבלה~\ref{tab:compet}) \\
אדם מחליט, המערכת מתעדת & הגנה משפטית ותדמיתית: לכל תמונה יש רשומת החלטה & תואם את אמנת פריז~\cite{paris2023} ואת סעיף \num{50}~\cite{euaiact2026} \\
תקן אתי והסמכה, בעברית ובערבית & יתרון מול קוראים ומפרסמים; אין ספק מקומי שמציע תקן אתי ותווית לקורא בעברית ובערבית & רק \num{12\%} מהציבור נוח עם חדשות שהופקו במכונה~\cite{reuters2025}; לאף מתחרה אין תקן (טבלה~\ref{tab:compet}) \\
\bottomrule
\end{tabularx}
\end{table}

טבלה~\ref{tab:exec} מרכזת את היעדים שהדוח מתחייב אליהם. לפיהם נבדוק בסוף הפרויקט אם הוא הצליח.


\begin{table}[!htb]\centering
\caption{היעדים המחייבים}\label{tab:exec}
\small
\begin{tabularx}{\linewidth}{@{}Y l@{}}
\toprule
\textbf{מדד} & \textbf{יעד} \\
\midrule
שיעור גילוי של תמונות מזויפות, במבחן על תמונות מהשטח & \num{90\%} לפחות \\
שיעור התרעות שווא על תמונות אמיתיות & \num{2\%} לכל היותר \\
זמן בדיקה אוטומטית (חציון) & עד \num{60} שניות \\
זמן החלטה כולל עם בודק אנושי (חציון) & עד \num{15} דקות \\
ירידה בפרסום תמונות מזויפות במערכות הפיילוט & \num{80\%} לפחות \\
משך הפרויקט עד גרסה מסחרית & \num{18} חודשים \\
תקציב, כולל רזרבה של \num{12\%} & \num{9.6} מיליון ₪ \\
\bottomrule
\end{tabularx}
\end{table}

\textbf{קנה מידה והבקשה.} היעד לשלוש שנים: \num{1,000} מערכות חדשות ו-\num{50} מיליון תמונות בחודש. אנו מבקשים לאשר את הדוח כבסיס חוזי; התוצרים וקריטריוני הקבלה בפרק~\ref{ch:deliv}, וטבלת העמידה בדרישות בנספח~\ref{app:compliance}.
